\documentclass[10pt,a4paper]{amsart}
\usepackage{amsmath,amssymb,amsthm,graphicx}
\usepackage{verbatim}
\newenvironment{reference}{\comment}{\endcomment}
%\newenvironment{reference}{}{}
\newenvironment{slogan}{\comment}{\endcomment}
\newenvironment{history}{\comment}{\endcomment}

% For commutative diagrams we use Xy-pic
\usepackage[all]{xy}

% We use 2cell for 2-commutative diagrams.
\xyoption{2cell}
% Theorem environments.
%
\theoremstyle{plain}
\newtheorem{theorem}[subsection]{Theorem}
\newtheorem{proposition}[subsection]{Proposition}
\newtheorem{lemma}[subsection]{Lemma}

\theoremstyle{definition}
\newtheorem{definition}[subsection]{Definition}
\newtheorem{example}[subsection]{Example}
\newtheorem{exercise}[subsection]{Exercise}
\newtheorem{situation}[subsection]{Situation}

\theoremstyle{remark}
\newtheorem{remark}[subsection]{Remark}
\newtheorem{remarks}[subsection]{Remarks}

\numberwithin{equation}{subsection}

% Macros
%
\def\lim{\mathop{\mathrm{lim}}\nolimits}
\def\colim{\mathop{\mathrm{colim}}\nolimits}
\def\Spec{\mathop{\mathrm{Spec}}}
\def\Hom{\mathop{\mathrm{Hom}}\nolimits}
\def\Ext{\mathop{\mathrm{Ext}}\nolimits}
\def\SheafHom{\mathop{\mathcal{H}\!\mathit{om}}\nolimits}
\def\SheafExt{\mathop{\mathcal{E}\!\mathit{xt}}\nolimits}
\def\Sch{\mathit{Sch}}
\def\Mor{\mathop{\mathrm{Mor}}\nolimits}
\def\Ob{\mathop{\mathrm{Ob}}\nolimits}
\def\Sh{\mathop{\mathit{Sh}}\nolimits}
\def\NL{\mathop{N\!L}\nolimits}
\def\CH{\mathop{\mathrm{CH}}\nolimits}
\def\proetale{{pro\text{-}\acute{e}tale}}
\def\etale{{\acute{e}tale}}
\def\QCoh{\mathit{QCoh}}
\def\Ker{\mathop{\mathrm{Ker}}}
\def\Im{\mathop{\mathrm{Im}}}
\def\Coker{\mathop{\mathrm{Coker}}}
\def\Coim{\mathop{\mathrm{Coim}}}

% Boxtimes
%
\DeclareMathSymbol{\boxtimes}{\mathbin}{AMSa}{"02}

%
% Macros for moduli stacks/spaces
%
\def\QCohstack{\mathcal{QC}\!\mathit{oh}}
\def\Cohstack{\mathcal{C}\!\mathit{oh}}
\def\Spacesstack{\mathcal{S}\!\mathit{paces}}
\def\Quotfunctor{\mathrm{Quot}}
\def\Hilbfunctor{\mathrm{Hilb}}
\def\Curvesstack{\mathcal{C}\!\mathit{urves}}
\def\Polarizedstack{\mathcal{P}\!\mathit{olarized}}
\def\Complexesstack{\mathcal{C}\!\mathit{omplexes}}
% \Pic is the operator that assigns to X its picard group, usage \Pic(X)
% \Picardstack_{X/B} denotes the Picard stack of X over B
% \Picardfunctor_{X/B} denotes the Picard functor of X over B
\def\Pic{\mathop{\mathrm{Pic}}\nolimits}
\def\Picardstack{\mathcal{P}\!\mathit{ic}}
\def\Picardfunctor{\mathrm{Pic}}
\def\Deformationcategory{\mathcal{D}\!\mathit{ef}}
\usepackage{cleveref}
\crefname{originalreference}{original reference}{original references}
\Crefname{originalreference}{Original reference}{Original references}
\makeatletter
\newlabel{lemma-finite-Galois}{{09I1}{source}{}{sourcepointer-1}{}}
\newlabel{lemma-finite-Galois@cref}{{[originalreference][0][]09I1}{[1][0][]source}{}{}{}}
\makeatother
\crefname{originalreference}{original reference}{original references}
\Crefname{originalreference}{Original reference}{Original references}
\makeatletter
\newlabel{lemma-finite-goes-up}{{09G5}{source}{}{sourcepointer-2}{}}
\newlabel{lemma-finite-goes-up@cref}{{[originalreference][0][]09G5}{[1][0][]source}{}{}{}}
\makeatother
\crefname{originalreference}{original reference}{original references}
\Crefname{originalreference}{Original reference}{Original references}
\makeatletter
\newlabel{lemma-galois-goes-up}{{09I2}{source}{}{sourcepointer-3}{}}
\newlabel{lemma-galois-goes-up@cref}{{[originalreference][0][]09I2}{[1][0][]source}{}{}{}}
\makeatother
\crefname{originalreference}{original reference}{original references}
\Crefname{originalreference}{Original reference}{Original references}
\makeatletter
\newlabel{lemma-galois-over-fixed-field}{{09I3}{source}{}{sourcepointer-4}{}}
\newlabel{lemma-galois-over-fixed-field@cref}{{[originalreference][0][]09I3}{[1][0][]source}{}{}{}}
\makeatother
\crefname{originalreference}{original reference}{original references}
\Crefname{originalreference}{Original reference}{Original references}
\makeatletter
\newlabel{lemma-lift-maps}{{0BME}{source}{}{sourcepointer-5}{}}
\newlabel{lemma-lift-maps@cref}{{[originalreference][0][]0BME}{[1][0][]source}{}{}{}}
\makeatother
\title{Fundamental theorem of finite Galois theory: subgroups correspond to subextensions}
\author{The Stacks Project Authors}
\date{stacks-project:fields.tex:local-snapshot-20261012; GNU FDL 1.2 or later}
\begin{document}
\maketitle
\noindent\textit{Editorial note.} All delivered mathematical text is the author's own, copied from the pinned Stacks Project chapter file fields.tex; no statement or proof step was rewritten, reordered or removed.\par\medskip
\noindent\textit{Editorial note.} The counted proof is the author's complete argument: Both squeeze ingredients are the two lemmas proved inside the span's own first paragraph: |Gal(L/M)| = [L : M] for subextensions and [L : L\textasciicircum{}H] = |H| for subgroups.\par\medskip
\noindent\textit{Editorial note.} The label theorem-galois-theory names the counted statement, and the proof environment immediately following it in the source is delivered verbatim as its proof.\par\medskip
\noindent\textit{Editorial note.} The counted statement is carried by the author's own theorem environment in the chapter source.\par\medskip
\noindent\textit{Allowed standard prerequisites (editorial).}\begin{itemize}
\item Context spans: Gal(E/F) notation and the fixed field L\textasciicircum{}H
\item Same chapter: lemma-galois-goes-up (L/M is Galois for every subextension L/M/K of a Galois L/K), lemma-finite-goes-up (finiteness goes up), lemma-finite-Galois (finite extension is Galois iff |Aut| = degree), lemma-galois-over-fixed-field ([L : L\textasciicircum{}H] = |H| for finite H), lemma-lift-maps (restriction of automorphisms to subextensions)
\end{itemize}\medskip
\ifdefined\hypertarget\hypertarget{sourcepointer-1}{}\fi
\noindent\textit{Source reference.} Original reference 09I1: sources/stacks-fields/pinned-original-source.tex, source lines 2881--2881. Prerequisite pointer invoked at 3019; it names a standard fact of the same chapter that the counted argument uses, and every step of the counted proof built on it is retained verbatim in the delivered body.\par\medskip
\ifdefined\hypertarget\hypertarget{sourcepointer-2}{}\fi
\noindent\textit{Source reference.} Original reference 09G5: sources/stacks-fields/pinned-original-source.tex, source lines 668--668. Prerequisite pointer invoked at 3018; it names a standard fact of the same chapter that the counted argument uses, and every step of the counted proof built on it is retained verbatim in the delivered body.\par\medskip
\ifdefined\hypertarget\hypertarget{sourcepointer-3}{}\fi
\noindent\textit{Source reference.} Original reference 09I2: sources/stacks-fields/pinned-original-source.tex, source lines 2911--2911. Prerequisite pointer invoked at 3016; it names a standard fact of the same chapter that the counted argument uses, and every step of the counted proof built on it is retained verbatim in the delivered body.\par\medskip
\ifdefined\hypertarget\hypertarget{sourcepointer-4}{}\fi
\noindent\textit{Source reference.} Original reference 09I3: sources/stacks-fields/pinned-original-source.tex, source lines 2944--2944. Prerequisite pointer invoked at 3021; it names a standard fact of the same chapter that the counted argument uses, and every step of the counted proof built on it is retained verbatim in the delivered body.\par\medskip
\ifdefined\hypertarget\hypertarget{sourcepointer-5}{}\fi
\noindent\textit{Source reference.} Original reference 0BME: sources/stacks-fields/pinned-original-source.tex, source lines 2163--2163. Prerequisite pointer invoked at 3034; it names a standard fact of the same chapter that the counted argument uses, and every step of the counted proof built on it is retained verbatim in the delivered body.\par\medskip
\begin{definition}
\label{definition-galois-group}
If $E/F$ is a Galois extension, then the group $\text{Aut}(E/F)$ is
called the {\it Galois group} and it is denoted $\text{Gal}(E/F)$.
\end{definition}

\noindent
Let $G$ be a group acting on a field $K$ (by field automorphisms).
We will often use the notation
$$
K^G = \{x \in K \mid \sigma(x) = x \ \forall \sigma \in G\}
$$
and we will call this the {\it fixed field} for the action of $G$ on $K$.

\begin{theorem}[Fundamental theorem of Galois theory]
\label{theorem-galois-theory}
Let $L/K$ be a finite Galois extension with Galois group $G$.
Then we have $K = L^G$ and the map
$$
\{\text{subgroups of }G\}
\longrightarrow
\{\text{subextensions }L/M/K\},\quad
H \longmapsto L^H
$$
is a bijection whose inverse maps $M$ to $\text{Gal}(L/M)$.
The normal subgroups $H$ of $G$ correspond exactly to those
subextensions $M$ with $M/K$ Galois.
\end{theorem}

\begin{proof}
By Lemma \ref{lemma-galois-goes-up} given a subextension $L/M/K$
the extension $L/M$ is Galois. Of course $L/M$ is also finite
(Lemma \ref{lemma-finite-goes-up}). Thus $|\text{Gal}(L/M)| = [L : M]$
by Lemma \ref{lemma-finite-Galois}.
Conversely, if $H \subset G$ is a finite subgroup, then
$[L : L^H] = |H|$ by Lemma \ref{lemma-galois-over-fixed-field}.
It follows formally from these two observations that we obtain
a bijective correspondence as in the theorem.

\medskip\noindent
If $H \subset G$ is normal, then $L^H$ is fixed by the action of
$G$ and we obtain a canonical map $G/H \to \text{Aut}(L^H/K)$.
This map has to be injective as $\text{Gal}(L/L^H) = H$. Hence
$|G/H| = [L^H : K]$ and $L^H$ is Galois by
Lemma \ref{lemma-finite-Galois}.

\medskip\noindent
Conversely, assume that $K \subset M \subset L$ with $M/K$ Galois.
By Lemma \ref{lemma-lift-maps} we see that every
element $\tau \in \text{Gal}(L/K)$ induces an element
$\tau|_M \in \text{Gal}(M/K)$. This induces a homomorphism
of Galois groups $\text{Gal}(L/K) \to \text{Gal}(M/K)$ whose
kernel is $H$. Thus $H$ is a normal subgroup.
\end{proof}

\end{document}
